\scanpage{241}
\mspage{121}
% l. 1
que \unsure{$\xi_0$} \unsure{correspondent} \inserted{$1-1$} aux
$\xi'\in H^1(S,\unsure{\mathcal H'_S})$.

% l. 2
Pour un tel \deleted{\illegible{1}} $\xi'$, soit
$\delta'(\xi')\in H^2(S,\mu)$
% l. 3
l'obstruction à le relever en $H^1(S,\unsure{\mathcal G'_S|\mathcal G_S})$,
% l. 4
et $\delta(\xi')=\delta(\xi)$ l'obstruction à le relever
% l. 5
en $H^1(S,\mathcal G_S)$. \unsure{Soit encore|C'est encore}
% l. 6
\[
\delta'(\xi')=\delta(\xi')-\rho_0
\]
% l. 7
i.e.
\[
\delta(\xi')=\delta(\xi)=\delta'(\xi')+\rho_0
\]
% l. 8
et on veut $\delta(\xi)=\rho$ i.e.
% l. 9
\[
\deleted{\unsure{\delta(\xi')}=}\quad
\delta'(\xi')=\rho-\rho_0
\]
% l. 10
et la question revient encore à \deleted{\illegible{3}} \inserted{la}
% l. 11
surjectivité de $H^1(S,\mathcal H'_S)\longrightarrow H^2(S,\unsure{\mu})$
\deleted{\illegible{1}}
% l. 12
--- je ne m'en tire pas. Il faudrait
% l. 13
consulter des gens compétents, comme
% l. 14
Giraud ou \unsure{Larry Breen}. On sent qu'il
% l. 15
faudrait travailler avec un \unsure{Gr-champ} $\unsure{\underline{\underline{\mathcal H}}}$
% l. 16
de coefficients, d'invariants $\underline{\pi}_0\simeq\Gamma$ et
$\underline{\pi}_1\simeq\mu$,
% l. 17
(mais pas nécessairement un champ de \unsure{Picard}!)
% l. 18
\deleted{\unsure{avec}} et
$H^1(S,\underline{\underline{\mathcal H}})\simeq\text{Classes d'applications}$
% l. 19
de $S$ dans $B\underline{\underline{\mathcal H}}$, qui est un espace
% l. 20
connexe avec $\pi_1(B\underline{\underline{\mathcal H}})\simeq\Gamma$,
$\pi_2(B\underline{\underline{\mathcal H}})=\mu$.
% l. 21
Ici la classe de \unsure{Postnikoff} dans $H^3(\Gamma,\mu)$ est nulle,
\margin{240}

\scanpage{242}
% l. 1
[mais peut-être n'y a-t-il pas lieu de
% l. 2
trop se raccrocher à cette \unsure{hypothèse},
% l. 3
correspondant à l'existence d'une \inserted{\unsure{ext.}} $\mathcal G$ de $\mathcal H$
\inserted{par $\pi_1(\mathcal H^0)$} \deleted{\illegible{1}}
% l. 4
qui redonne l'extension universelle de
% l. 5
$\mathcal H^0$ par $\pi_1(\mathcal H^0)$] on a un \unsure{hom.}
% l. 6
$B\mathcal H\longrightarrow B\underline{\underline{\mathcal H}}$
\deleted{\illegible{1}} qui induit un iso pour
% l. 7
les $\pi_i$ pour $i\leq2$, et pour $i\leq n$
(\unsure{où} $n\geq2$ \inserted{\unsure{est donné}})
% l. 8
\unsure{sss|ssi} $\pi_i(B\mathcal H)=0$ pour $2<i\leq n$, i.e.
% l. 9
$\pi_i(\mathcal H)=0$ pour $2\leq i\leq n-1$ [dans le cas
% l. 10
qui nous intéresse $\mathcal H=A_0$, on peut
% l. 11
prendre $n=3$). Ceci implique que pour
% l. 12
tout C-W-complexe $S$ de dim $\leq n$,
% l. 13
on a \inserted{(Hot désignant les morphismes dans
la cat. homotopique \unsure{non ponctuée})}
% l. 14
\[
\mathrm{Hot}(S,B\mathcal H)\xrightarrow{\sim}
\mathrm{Hot}(S,B\underline{\underline{\mathcal H}})
\]
% l. 15
\unsure{c.à.d.|i.e.}
\[
H^1(S,\mathcal H)\xrightarrow{\simeq}
H^1(S,\underline{\underline{\mathcal H}})
\]

% l. 16
Mais si l'invariant de \unsure{Postnikoff}, donc
% l. 17
$\unsure{K|\kappa}\in H^3(\Gamma,\mu)$ est nul,
\deleted{\illegible{1}} (ainsi que l'action
% l. 18
de $\Gamma$ sur $\mu$) alors \unsure{sauf erreur}
$B\underline{\underline{\mathcal H}}$ s'iden%
% l. 19
tifie à un produit $K(\Gamma,1)\times K(\unsure{\pi|\mu},2)$,
\deleted{\unsure{et}}
% l. 20
(cette ``identification'' dépend \unsure{justement} du
\margin{241}

\scanpage{243}
\mspage{122}
% l. 1
choix \unsure{dans|des} $H^2(\pi_1,\pi_2)\simeq H^2(\Gamma,\mu)$ !), et par
% l. 2
suite
\[
H^1(S,\underline{\underline{\mathcal H}})\xrightarrow{\simeq}
H^1(S,\Gamma)\times H^2(S,\mu)
\]
% l. 3
pour tt espace $S$, donc pour $S$ un C-W-%
% l. 4
complexe de dim $\leq n$
% l. 5
\[
\unsure{(7_1)|(\dagger_1)}\qquad
H^1(S,\mathcal H)\xrightarrow{\simeq}H^1(S,\Gamma)\times H^2(S,\mu)
\]

% l. 6
Donc la classification des fibrés en \unsure{sphères}
% l. 7
topologiques sur un espace topologique $S$,
% l. 8
pour un C-W-complexe $S$ de dim $\leq3$,
% l. 9
marche bel et bien. Bien entendu, le
% l. 10
fait qu'on soit obligé ici de se borner à
% l. 11
\deleted{\illegible{2}} $S$ de dim $\leq3$ tient au fait
% l. 12
que nous n'avons trouvé \deleted{\unsure{qu'}}
(via $\underline{\underline{\mathcal H}}$)
% l. 13
qu'une approximation $2$-\unsure{isotopique|homotopique} (et non
% l. 14
$\infty$-\unsure{isotopique|homotopique}) de la catégorie des $2$-sphères
% l. 15
topologiques.

% l. 16
Je voudrais reprendre la classification
% l. 17
pour un CW-complexe $S$ quelconque, en
% l. 18
utilisant le fait que l'inclusion
% l. 19
\[
\widetilde{\mathrm{GP}}(1,\mathbb C)\hookrightarrow A_0
\]
\margin{242}

\scanpage{244}
% l. 1
est une équivalence d'homotopie, et induit
% l. 2
donc une bijection
% l. 3
\[
H^1(S,\widetilde{\mathrm{GP}}(1,\mathbb C))\longrightarrow H^1(S,A_0)
\]
% l. 4
et $\unsure{\overline{\mathrm m}|\simeq}$ une $\infty$-équivalence des $\infty$-catégories des
% l. 5
\deleted{fibrés} \inserted{torseurs} sur $S$, de groupes
$\widetilde{\mathrm{GP}}$ (correspondant
% l. 6
aux fibrés en sphères conformes, ou en
% l. 7
droites projectives sur une \deleted{corps} $\mathbb R$-algèbre
\inserted{\unsure{en principe}}
% l. 8
isomorphe à $\mathbb C$) et de groupe $A_0$ --- correspondant
% l. 9
aux fibrés en sphères sur $S$.

% l. 10
Le premier invariant d'un fibré de groupe $\widetilde{\mathrm{GP}}$
% l. 11
est un \deleted{fibré}
\unsure{\inserted{torseur $T$ sous $(\mathbb Z/2\mathbb Z)_S$
\deleted{\illegible{2}} \unsure{défini} (\unsure{à isom.} près) par un}}
% l. 12
\uline{\unsure{invariant numérique}} $\chi\ (\in H^1(S,\mathbb Z/2\mathbb Z))$.
\margin{i.e. un \unsure{rev. étale|rev. double} de degré $2$.
La donnée de $\chi$ revient à la donnée d'un système d'orientations
\unsure{tordues} sur $S$.}
% l. 13
\inserted{Ce torseur $T$} servira à tordre $\mathbb C$
(via l'opération \inserted{fidèle} de $\mathbb Z/2\mathbb Z$ sur
% l. 14
la $\mathbb R$-\unsure{algèbre} $\mathbb C$), \deleted{et} d'où un fibré loc.
\deleted{\illegible{1}}
% l. 15
constant en \deleted{corps} extensions \inserted{quadratiques}
de $\mathbb R$ \deleted{\illegible{1}},
% l. 16
soit $C$, et on \deleted{prendra le} \unsure{enrichira|enrichit} $S$
% l. 17
par le faisceau \inserted{$\underline C_S$ (ou $\underline C$)} des sections continues de
% l. 18
$C$ (ce qui est une façon de tordre
% l. 19
\deleted{\unsure{$\mathrm{SL}(2,$}} $\mathrm{GP}(1,\mathbb C)_S$
en $\mathrm{GP}(1,C)_S=\mathrm{GP}(1,\underline C)$))
% l. 20
et il s'agit de \deleted{\unsure{donc}} décrire de façon
% l. 21
compréhensible \deleted{la catégorie} \illegible{2} \unsure{aux bons}
% l. 22
\deleted{\illegible{1}} \unsure{$n$-isotopiques}, pour $n=\ldots$ --- la catégorie
\margin{243}

\scanpage{245}
\mspage{123}
% l. 1
--- dépendant du choix de $\chi$ ou $\underline C$ --- des
% l. 2
algèbres d'Azumaya de rang $4$ sur $\underline C$,
% l. 3
ou encore des fibrés en droites projectives
% l. 4
sur $\underline C$, ou des torseurs sous $\mathrm{GP}(1,\underline C)$.

% l. 5
Or utilisant \deleted{$\mathrm{SL}(2,\underline C)\longrightarrow
\mathrm{GP}(1,\underline C)$} \inserted{la suite exacte}
% l. 6
\[
1\longrightarrow\mu_{2S}\longrightarrow\mathrm{SL}(2,\underline C_S)
\longrightarrow\mathrm{GP}(1,\underline C_S)\longrightarrow1
\]
% l. 7
on associe à un objet de la catégorie
% l. 8
--- disons un torseur sous $\mathrm{GP}(1,\underline C_S)$ --- la
% l. 9
catégorie \inserted{fibrée (sur des ouverts variables de $S$)} de ses relèvements à
$\mathrm{SL}(2,\underline C_S)$,
% l. 10
qui est une \uline{gerbe} \inserted{$\mathcal G$} liée par $\unsure{\mu_2|\mu}$.
% l. 11
C'est cette gerbe qui est le deuxième
% l. 12
invariant complet --- plus fin que
% l. 13
sa classe \deleted{dans $H^2($} \inserted{d'équivalence, qui s'identifie à un}
% l. 14
$\xi\in H^2(S,\mu_2)=H^2(S,\{\pm1\})$
(jouant le rôle
% l. 15
d'un groupe de Brauer)

% l. 16
Un isomorphisme de fibrés, d'invariants
% l. 17
$(T,\mathcal G)$, $(T',\mathcal G')$, définira un isom $T\simeq T'$
% l. 18
(\unsure{c.à.d.} un isom \deleted{$\mathrm{GP}(1,C$}
$\underline C_T\simeq\underline C_{T'}$ d'où
% l. 19
un isom $\mathrm{GP}(1,\underline C_T)\simeq\mathrm{GP}(1,\underline C_{T'})$) --- et
% l. 20
une équivalence de gerbes --- il \unsure{faudrait|faudra}
\margin{244}

\scanpage{246}
% l. 1
expliciter que pour deux isomorphismes
% l. 2
$f,g$ de torseurs, \deleted{et} d'où
% l. 3
\diagram{[À gauche : $f_T,g_T:T\longrightarrow T'$, avec
\unsure{deux flèches parallèles} dirigées vers $T'$, des signes
$\simeq$ au-dessus et plusieurs reprises de traits.
À droite : $f_{\mathcal G},g_{\mathcal G}:\mathcal G\longrightarrow\mathcal G'$,
avec deux flèches parallèles dirigées vers $\mathcal G'$ et des signes $\simeq$.]}
% l. 4
et pour tte homotopie $h_t$ $(0\leq t\leq1)$ de
% l. 5
$f$ à $g$, on trouve $f_T=g_T$ et un
% l. 6
isom $h_*:f_{\mathcal G}\xrightarrow{\sim}g_{\mathcal G}$ d'équivalences de
% l. 7
gerbes, qui ne dépend que de la classe
% l. 8
d'homotopie de cette homotopie. On
% l. 9
\unsure{trouve|trouvera} ainsi \unsure{un|une} $2$-foncteur, de la
% l. 10
$2$-catégorie des \deleted{fibrés} torseurs sur $S$
% l. 11
de groupe $\widetilde{\mathrm{GP}}(1,\mathbb C)$, vers la $2$-catégorie
% l. 12
formée des couples \inserted{$(T,\mathcal G)$} d'un
\deleted{torseur} \inserted{système} d'%
% l. 13
orientations tordues, \deleted{et} ($\Longleftrightarrow$ d'un torseur sur
% l. 14
$S$ de groupe \deleted{\illegible{1}} \inserted{$\mathbb Z/2$}) et d'une
% l. 15
\deleted{\unsure{$\{\pm1\}=$}} $\mu_2(\mathbb C)$-gerbe $\mathcal G$ sur $S$. Ce
% l. 16
$2$-foncteur, sauf erreur, est $2$-fidèle
% l. 17
sous les conditions \unsure{explicitées} plus
% l. 18
haut (\deleted{\illegible{2}} (dim $S\leq4$), et
% l. 19
est $3$-fidèle (i.e. l'application injective
% l. 20
$H^1(S,\widetilde{\mathrm{GP}})\unsure{\longrightarrow|{-}}
H^1(S,\mathbb Z/2\mathbb Z)\times H^2(S,\mu_2)$ est surjective)
\margin{245}

\scanpage{247}
\mspage{124}
% l. 1
si on a \unsure{même} dim $S\leq3$. On s'attend \unsure{qu'il|qu'ils}
% l. 2
soit $1$-fidèle dès que dim $S\leq5$, $0$-fidèle
% l. 3
dès que dim $S\leq6$.
\margin{??}
% A vertical marginal stroke accompanies approximately lines 1--3.
[\uline{NB} \inserted{l'assertion} ``$0$-fidèle''
\inserted{\unsure{ci-dessus|ci-dessous}} signifie
% l. 4
que si on a deux homotopies $h,h'$ de $f$ à $g$,
% l. 5
telles que $h_*=h'_*,$ alors $h$ et $h'$ sont homotopes.
% l. 6
La condition ``$1$-fidèle'' signifie \inserted{$0$-fidèle et} que tout
% l. 7
isom de $f_{\mathcal G}$ à $g_{\mathcal G}$
\deleted{\unsure{provient d'un}} est de la
% l. 8
forme $h_*$ (avec $h$ bien déterminé à hom. près,
% l. 9
par la condition précédente de $0$-fidélité). La
% l. 10
condition $2$-fidèle signifie que \deleted{\unsure{toute équivalence}}
\inserted{de plus, pour}
% l. 11
\deleted{\unsure{donnés}} $\varphi:T\simeq T'$,
$\unsure{\psi_{\mathcal G}|\psi}:\mathcal G\xrightarrow{\simeq}\mathcal G'$,
il existe
% l. 12
un isom de fibrés $f$ \unsure{tel} que $\varphi=f_T$,
\deleted{\illegible{1}}
% l. 13
et un isom $\psi\simeq f_{\mathcal G}$. Enfin $3$-fidèle signifie
% l. 14
que de plus, pour tout couple $T,\mathcal G$,
% l. 15
$\exists\ T',\mathcal G'$ provenant d'un fibré, \deleted{et de}
% l. 16
et un isom $T\simeq T'$ et une équivalence
% l. 17
$\mathcal G\simeq\mathcal G'$. Ces notions catégoriques
% l. 18
doivent \unsure{s'interpréter|interpréter} simplement que, si
% l. 19
on regarde $\mathcal H=\widetilde{\mathrm{GP}}(1,\mathbb C)
\xrightarrow{\alpha}\underline{\underline{\mathcal H}}$
--- où $\underline{\underline{\mathcal H}}$
% l. 20
est déduit de $\mathcal H$ en ``tuant les \unsure{g\textsuperscript{o}-pes}
% l. 21
d'homotopie $\pi_i$ pour $i\geq2$'', de sorte que
\margin{246}

\scanpage{248}
% l. 1
$\alpha$ induit un isom des $\pi_i$ pour $i\leq1$
% l. 2
(et $\unsure{\overline{\mathrm m}|=}$ pour $i\leq n$ si
$\pi_i(\mathcal H)=0$ pour $1<i\leq n$)
% l. 3
\deleted{et $B_\alpha$ un} d'où
% l. 4
\[
B_\alpha:B\mathcal H\longrightarrow B\underline{\underline{\mathcal H}}
\]
\margin{\uline{NB} $B\underline{\underline{\mathcal H}}$ est déduit de
$B\mathcal H$ en tuant les $\pi_i$ pour $i\geq3$.}
% l. 5
(induisant un isom des $\pi_i$ pour $i\leq n+1$),
% l. 6
alors prenant pour \deleted{tout} \inserted{un espace donné} $S$ l'application
% l. 7
correspondante des espaces d'applications continues
% l. 8
\[
\deleted{\mathrm{Hot}}
\underline{\mathrm{Cont}}(S,B\mathcal H)
\unsure{\longrightarrow|{-}}
\underline{\mathrm{Cont}}(S,B\underline{\underline{\mathcal H}})
\]
% l. 9
celle ci induit un isom pour les
% l. 10
$\pi_i$ $(0\leq i\leq2)$. \deleted{\unsure{lemme}} Il semblerait que
% l. 11
$0$-fidélité est une condition d'injectivité
% l. 12
pour les $\pi_2$, $1$-fidèle la bijectivité pour les
% l. 13
$\pi_2$ et l'injectivité pour les $\pi_1$, $2$-fidèle
% l. 14
la bijectivité pour $\pi_2$ et $\pi_1$ et l'injectivité
% l. 15
pour $\pi_0$, $3$-fidèle la bijectivité pour $\pi_0,\pi_1,\pi_2$.

% l. 16
Comme $\pi_i(\mathrm{Cont}(S,B))\simeq
\unsure{\mathrm{Ho}|\mathrm{Hot}}(S\times S^i,B)$,
% l. 17
\deleted{on voit} \inserted{il semblerait}
qu'on est conduit, pour la $0$-%
% l. 18
fidélité, de faire l'hypothèse draco%
% l. 19
nienne que $S\times S^2$ de dim $\leq4$, pour la
% l. 20
$1$-fidélité que $S\times S^2$ de dim $\leq3$ i.e. $S$
% l. 21
de dim $\leq1$ ? (\unsure{ceci|cela} impliquerait alors
% l. 22
la $3$-fidélité\ldots)
\margin{247}

\scanpage{249}
\mspage{125}
% l. 1
Prenons l'analogue arithmétique, d'une description
% l. 2
(plus ou moins ``fidèle'') \deleted{par \illegible{1}}
de nature ``\unsure{profinie}''
% l. 3
des droites projectives \deleted{définies} (éventuellement
% l. 4
tordues) définies sur un corps de type
% l. 5
fini $K$ \inserted{sur $\mathbb Q$ (plus gén. sur un schéma $S$ quelconque)},
celles ci forment a priori une catégorie
% l. 6
sans plus --- un groupoïde (pas nécessairement
% l. 7
connexe) je ne sais pas en faire une
% l. 8
$2$-catégorie raisonnablement, car pour deux
% l. 9
isom $f,g$ \inserted{\unsure{$(=f\circ u)$}} de tels torseurs, définir les
% l. 10
homotopies de $f$ à $g$, i.e. pour une ``forme''
% l. 11
$G$ de $\mathrm{GP}(1)_S$ et une section $u$ de $G$,
% l. 12
définir les ``homotopies'' de $u$ à l'%
% l. 13
identité, \deleted{et les classes \inserted{d'homotopie} de telles \illegible{1}}
\inserted{ou plutôt} une
% l. 14
notion qui remplace la notion de classe d'%
% l. 15
homotopie de telles homotopies. \deleted{Mais} \inserted{Ce serait}
% l. 16
sans doute un relèvement de $u$ en une
\deleted{
% l. 17
\unsure{Si} $S$ est \illegible{1} d'un pt géométrique
% l. 18
\illegible{1} \unsure{tte section} définit une \illegible{1}
% l. 19
\unsure{si $K$ est alg.} Pour un corps de base
% l. 20
alg. clos \unsure{tout au moins} (les \unsure{sections}
% l. 21
définissant un pt géométrique) on
% l. 22
peut songer à prendre les ``classes de
% l. 23
chemins de $f$ à \illegible{1} de $u$ à
}
% l. 24
section du \deleted{\unsure{$\widetilde{\mathrm{GP}}$}}
\unsure{revêtement} simplement
% l. 25
connexe $\widetilde G$ de $G$ !
\margin{248}

\scanpage{250}
% l. 1
je tiens là quelque chose d'assez amusant,
% l. 2
mais que je ne vais pas poursuivre --- de
% l. 3
\unsure{toutes façons|toute façon} il est clair que la ``description''
% l. 4
des droites projectives (\deleted{\unsure{sur}} un corps de
% l. 5
type fini disons) à laquelle on \unsure{aboutit},
% l. 6
n'a rien de fidèle --- pas même $0$-fidèle !
% l. 7
Ainsi tous les autom de $\mathbb P^1_K$ provenant de
% l. 8
$\mathrm{SL}(2,K)$ seront identifiés à l'identité ---
% l. 9
c'est un peu brutal ! Mais je me
% l. 10
rends compte que le travail conceptuel
% l. 11
autour des thèmes ``Brauer'' n'est pas
% l. 12
terminé --- qu'il y a à comprendre
% l. 13
des choses, pour l'étude de la catégorie
% l. 14
(\inserted{qui devrait être} une \uuline{2}-catégorie !)
des Algèbres d'Azumaya
% l. 15
\deleted{\unsure{de}} degré \uline{fixé} (ici $4$)\ldots
\margin{249}
